%=============================================================================!
% 12.2  MANY STATES, ONE SYSTEM                                               !
%=============================================================================!
\section{Many states, one system}
\label{sec:sm-ergodic}

A simulation follows one system along one trajectory, and its averages
are averages over time. Statistical mechanics computes averages over the
states the system could be in, weighted by their probabilities. This
section sets out the two kinds of average and the assumption that they
agree, and then shows systems for which they do not.

\subsection*{Microstates, macrostates and ensembles}

The state of $N$ atoms is a point $\Gamma = (\vb{r}^N, \vb{p}^N)$ of phase
space (\cref{sec:ha-equations}), called a microstate. What is measured of
a macroscopic sample, its energy, volume and number of atoms, or its
temperature and pressure, fixes a macrostate, which an enormous number of
microstates share. An ensemble is the cloud of copies of \cref{sec:ha-liouville}:
many imagined copies of the system, all in the same macrostate and spread
over its microstates, with a probability density $\mathcal{P}(\Gamma)$ in phase
space: a small box of phase space, of volume $\dd\Gamma$, the product of
its $6N$ sides, holds the fraction $\mathcal{P}(\Gamma)\,\dd\Gamma$ of the
copies. The ensemble average of a quantity $A(\Gamma)$ is
\begin{equation*}
   \ens{A} = \int A(\Gamma)\,\mathcal{P}(\Gamma)\,\dd\Gamma ,
\end{equation*}
an integral over all $6N$ coordinates and momenta, the limit of a sum
over such boxes, as for the volume integrals of \cref{sec:md-ewald}, and
the time average
along one trajectory is
\begin{equation*}
   \overline{A} = \lim_{t_{\mathrm f}\to\infty}\frac{1}{t_{\mathrm f}}
   \int_0^{t_{\mathrm f}}A\bigl(\Gamma(t)\bigr)\,\dd t .
\end{equation*}
\Cref{sec:ha-liouville} showed that a cloud whose density depends on the
state only through its energy, $\mathcal{P} = f(H)$, does not change as the
copies move. The ensembles of statistical mechanics are such clouds, and
the hypothesis that makes them useful to a simulation is the ergodic
hypothesis: that a single trajectory, followed long enough, passes near
every microstate of its energy, spending equal times near equal volumes of
phase space at that energy, so that $\overline{A}$ equals $\ens{A}$ for
the ensemble spread evenly over those states (\cref{sec:sm-micro}). It can be tested, and it fails for some systems.

\subsection*{A chain that never shares}

$N$ atoms of unit mass sit in a line between two fixed walls, each joined
to its neighbours by springs of unit stiffness. With $q_j$ the displacement
of atom $j$ and $q_0 = q_{N+1} = 0$ at the walls, the spring on the left
of atom $j$ is stretched by $q_j - q_{j-1}$ and pulls it back with
$-(q_j - q_{j-1})$, and the spring on its right is stretched by $q_{j+1} -
q_j$ and pulls it on with $q_{j+1} - q_j$, as for the carts of
\cref{sec:os-modes}; together they pull it with the force $q_{j-1} - 2q_j
+ q_{j+1}$. A normal mode has every atom moving as $u_j\cos(\omega t +
\phi)$, with the acceleration $-\omega^2u_j\cos(\omega t + \phi)$, which
for unit mass must equal the force; cancelling the common factor
$\cos(\omega t + \phi)$, this needs
\begin{equation*}
   -\omega^2u_j = u_{j-1} - 2u_j + u_{j+1} .
\end{equation*}
With $u_j = \sin j\theta$, the addition formulae (\cref{eq:mo-addition})
give $u_{j-1} + u_{j+1} = \sin(j\theta - \theta) + \sin(j\theta + \theta)
= 2\sin j\theta\cos\theta$, so the condition holds with
\begin{equation*}
   \omega^2 = 2 - 2\cos\theta = 4\sin^2\tfrac12\theta ,
\end{equation*}
using $1 - \cos\theta = 2\sin^2\frac12\theta$ (\cref{ex:mo-identity}).
The wall at $j = 0$ is met, since $\sin0 = 0$, and the one at $j = N + 1$
needs $\sin(N + 1)\theta = 0$, so $\theta = k\pi/(N + 1)$ for a whole
number $k$. The values $k = 1$ to $N$ give $N$ different patterns, as many
as the chain has modes (\cref{sec:os-modes}); $k = N + 1$ gives $\sin
j\pi = 0$ for every atom, no motion at all, and larger $k$ give the same
patterns again, up to sign. The modes have the angular frequencies
$\omega_k = 2\sin(k\pi/(2N + 2))$, the positive root, since $k\pi/(2N +
2)$ lies between 0 and $\pi/2$. For $N =
32$, the slowest repeats every 66.0 units of time and the fastest every
3.15. As \cref{sec:os-modes} showed, every motion of the harmonic chain
is a sum of normal modes, each keeping its own amplitude, so a chain
started with its energy in one mode keeps it there for ever: two starts
with the same energy $E$, one in the first mode and one in the second,
give the energy in the first mode the time averages $E$ and 0. The
harmonic chain is not ergodic.

Real bonds are not perfect springs. Adding an energy $\frac13\alpha x^3$
to each spring stretched by $x$ couples the modes, and the question is
whether the coupling spreads the energy among them. Fermi, Pasta and Ulam
asked it in what Dauxois calls the first numerical experiment, programmed
and run by Mary Tsingou\cite{Dauxois2008}. \Cref{fig:sm-ergodic}(a) repeats
it with \texttt{mdlab.chain}, for 32 atoms with $\alpha = 0.25$ and the
energy 0.0724 put into the first mode, followed by velocity Verlet with
steps of 0.05, over which the total energy varies by no more than
$\num{4.9e-5}$ of itself. The first mode's share falls below 0.1 by $t = 3950$, the
energy passing to modes 2, 3 and 4 and back, and by $t = 10\,225$ it has
returned to 0.981. Over the whole run, to $t = 20\,000$, about 300
periods of the first mode, modes 5 to 32 never hold more than 0.198 of the
energy; spread equally, they would hold $28/32 = 0.875$. A trajectory
started in this special state keeps its energy in the few slowest modes,
passing it among them and back to the first, for the whole run; a run a
hundred times longer, to $t = \num{2e6}$, still leaves only 0.086 of it on
average in modes 5 to 32, and never more than 0.220.

\begin{figure}[htb]
\centering
\includegraphics{ch12_ensembles/ergodic}
\caption{Time averages. (a) A chain of 32 atoms started with its energy
in the first normal mode: the share of the energy in modes 1 to 4 against
time, with the coupling $\alpha = 0.25$ (the four labelled curves), and for
the harmonic chain, whose first mode keeps all of it (grey dashed). (b) Liquid argon, 256
atoms: the running time average of the kinetic energy of three single
atoms, divided by the average over all atoms and times (dotted at 1).}
\label{fig:sm-ergodic}
\end{figure}

\subsection*{Time averages in a liquid}

The liquid provides a contrasting test over a finite observation time. \Cref{fig:sm-ergodic}(b) follows 256 atoms
of liquid argon, the model of \cref{sec:md-loop} at the density
$0.8\sigma^{-3}$ and about \SI{135}{\kelvin} (a temperature in the sense of
\cref{sec:sm-equipartition}), prepared as in
\cref{sec:cs-water}, for \SI{50}{\pico\second} at fixed energy. The
running time average of one atom's kinetic energy wanders at first and
then settles towards the average over all atoms and saved times. If
the atoms sample the same stationary distribution, this pooled average
estimates the ensemble average: after
\SI{50}{\pico\second} the time averages of the 256 atoms lie between 0.807
and 1.249 of it, with a standard deviation of 0.076. They are still
settling, and how long a run must be for an average to be trusted is the
question of Chapter~15. A crystal at a low temperature, whose
vibrations are nearly harmonic, is closer to the chain than to the
liquid, and sharing in it can be correspondingly slow.

\begin{practice}
A run that starts in a special state, such as a crystal on its lattice or
energy given to one motion only, may keep the memory of that start. The
start should be one the system could reach from many others, and its
averages should not depend on which of several such starts is used.
\end{practice}

\begin{keyidea}
An ensemble is a cloud of copies of a system with a density $\mathcal{P}(\Gamma)$
in phase space, and $\ens{A}$ its average; a simulation measures the time
average $\overline{A}$. The ergodic hypothesis, that the two agree, is
consistent with liquid argon over \SI{50}{\pico\second}; it fails for the
harmonic chain, whose modes never share energy. The chain of Fermi,
Pasta, Ulam and Tsingou also fails to share its energy equally over the
observation times tested here, returning much of it to the first mode;
that finite-time result does not establish its behaviour at infinite time.
\end{keyidea}

\notebook{12}{run the chain with any coupling and watch the modes share or not}

\begin{selfcheck}
Why is the period of the fastest mode of a long chain close to $\pi$?
\end{selfcheck}
